
Because $U_6$ has eigenphases $e^{in\pi/3}$, the regularized transport spectrum is
\begin{equation}
\label{eq:transport-spectrum}
\lambda_{y,n}^{-1}
=
y^2+\delta_{\mathrm{ph}}^2-2y\delta_{\mathrm{ph}}\cos\!\left(\frac{n\pi}{3}\right)+\varepsilon^2,
\qquad
n=0,\dots,5,
\end{equation}
for the eigenvalues $\lambda_{y,n}$ of $\mathcal Y_y^{(\varepsilon)}$. In particular,
\begin{lemma}[Finite admissible hexagon gap]
\label{lem:finite-admissible-hexagon-gap}
On the admissible cusp set
\[
y\in\left\{1,\frac13,\frac23\right\},
\]
the smallest singular value of the unregulated hexagon denominator is attained at $y=\frac23$ and $n=0$. Equivalently,
\[
\epsilon_f
:=
\min_{\substack{y\in\{1,\frac13,\frac23\}\\ n=0,\dots,5}}
\left|y-\delta_{\mathrm{ph}}e^{in\pi/3}\right|
=
\frac23-\delta_{\mathrm{ph}}.
\]
Hence the finite admissible hexagon has a strictly positive spectral gap $\epsilon_f$.
\end{lemma}

\begin{proof}
For $\varepsilon=0$, the eigenvalues in \Cref{eq:transport-spectrum} reduce to
$|y-\delta_{\mathrm{ph}}e^{in\pi/3}|^{-2}$. By \Cref{thm:transport-pole-algebraic},
$\delta_{\mathrm{ph}}$ is the unique admissible root in $(1/3,2/3)$. Hence the positive-real cusp
$y=2/3$ with $n=0$ is the nearest admissible point, while every other cusp or nonzero phase has
strictly larger distance. Therefore
\[
\epsilon_f=\frac23-\delta_{\mathrm{ph}},
\]
which is positive because $\delta_{\mathrm{ph}}<2/3$.
\end{proof}

\begin{definition}[Gap-stable admissibility]
\label{def:gap-stable-admissibility}
Let
\[
D_0:=D_{2/3}.
\]
We say that the admissible transport sector is gap-stable if the physical transport operator has the form
\[
D_f=D_0+V_{\mathrm{adm}},
\qquad
[P_{\mathrm{str}},V_{\mathrm{adm}}]=0,
\qquad
\|V_{\mathrm{adm}}\|<\frac{\epsilon_f}{2}.
\]
\end{definition}

\begin{proposition}[Structural gap lift]
\label{prop:gap-lift-full-operator}
If admissibility is gap-stable in the sense of \Cref{def:gap-stable-admissibility}, then for every
\[
\psi\in\mathrm{Ran}(P_{\mathrm{str}})
\]
one has
\[
\|D_f\psi\|
\ge
\frac{\epsilon_f}{2}\|\psi\|.
\]
Consequently
\begin{equation}
\label{eq:gap-lift-full-operator}
P_{\mathrm{str}} D_f^\dagger D_f P_{\mathrm{str}}
\ge
\frac{\epsilon_f^2}{4}\,P_{\mathrm{str}}
>
0.
\end{equation}
In particular, the final admissible restriction to $\mathrm{Ran}(\Padm)$ is strictly positive.
\end{proposition}

\begin{proof}
Since $\psi\in\mathrm{Ran}(P_{\mathrm{str}})$ and $[P_{\mathrm{str}},V_{\mathrm{adm}}]=0$,
\[
\|D_f\psi\|
\ge
\|D_0\psi\|-\|V_{\mathrm{adm}}\psi\|
\ge
\bigl(\epsilon_f-\|V_{\mathrm{adm}}\|\bigr)\|\psi\|
\ge
\frac{\epsilon_f}{2}\|\psi\|.
\]
Squaring gives \Cref{eq:gap-lift-full-operator}.
\end{proof}

\begin{equation}
\label{eq:mobius-quotient}
\frac{\lambda_{y,0}}{\lambda_{y,3}}
=
\frac{(y+\delta_{\mathrm{ph}})^2+\varepsilon^2}{(y-\delta_{\mathrm{ph}})^2+\varepsilon^2}.
\end{equation}
\[
\lim_{\varepsilon\to 0}\frac{\lambda_{y,0}}{\lambda_{y,3}}
=
\left(\frac{y+\delta_{\mathrm{ph}}}{y-\delta_{\mathrm{ph}}}\right)^2
\]
\begin{remark}[Near-pole hierarchy mechanism]
On the closed transport branch, the exact pole is the root of
\Cref{prop:closed-form-transport-pole}; the later comparison ledgers retain that same pole and
regularize only the finite Green kernel at $\varepsilon>0$.
Hence the Möbius quotient contains a built-in near-pole whenever $y$ approaches $\delta_{\mathrm{ph}}$, making large hierarchy ratios structurally plausible. In the closed transport grammar this is the hierarchy mechanism, while the bounded holonomy factors keep it from degenerating into an arbitrary tuning surface.
\end{remark}

\begin{proposition}[Finite positive transport algebra]
\label{prop:finite-positive-transport-algebra}
Let
\[
C:=\frac12\bigl(U_6+U_6^\dagger\bigr),
\qquad
a:=y^2+\delta_{\mathrm{ph}}^2+\varepsilon^2,
\qquad
b:=y\delta_{\mathrm{ph}}.
\]
Then
\[
D_y^\dagger D_y+\varepsilon^2=a\,\mathbf 1-2b\,C.
\]
The cosine operator $C$ has spectrum
\[
\mathrm{Spec}(C)=\left\{1,\frac12,-\frac12,-1\right\}
\]
and satisfies
\begin{equation}
\label{eq:c6-cosine-polynomial}
4C^4-5C^2+\mathbf 1=0.
\end{equation}
Consequently the regularized positive transport kernel reduces to the cubic polynomial
\begin{equation}
\label{eq:finite-positive-transport-kernel}
\mathcal Y_y^{(\varepsilon)}
=
\frac{
a^3-5ab^2
+2b(a^2-5b^2)C
+4ab^2C^2
+8b^3C^3
}{
(a^2-b^2)(a^2-4b^2)
}.
\end{equation}
\end{proposition}

\begin{proof}
Because $U_6$ is unitary,
\[
D_y^\dagger D_y
=
\bigl(y\,\mathbf 1-\delta_{\mathrm{ph}}U_6^\dagger\bigr)
\bigl(y\,\mathbf 1-\delta_{\mathrm{ph}}U_6\bigr)
=
\bigl(y^2+\delta_{\mathrm{ph}}^2\bigr)\mathbf 1
-y\delta_{\mathrm{ph}}(U_6+U_6^\dagger),
\]
which gives the first identity after adding $\varepsilon^2$. The eigenvalues of $C$ are $\cos(n\pi/3)$ for $n=0,\dots,5$, namely $1,\frac12,-\frac12,-1$. Hence $C$ obeys the minimal polynomial \Cref{eq:c6-cosine-polynomial}. The inverse of $a-2bc$ on this four-point spectrum is therefore represented by a unique cubic interpolation polynomial in $c$, and evaluating that interpolation yields \Cref{eq:finite-positive-transport-kernel}.
\end{proof}

\begin{proposition}[Strict positivity of the finite admissible $C_6$ Green kernel]
\label{prop:strict-positive-c6-green}
For each admissible cusp
\[
y\in\left\{1,\frac13,\frac23\right\}
\]
and each $\varepsilon\ge 0$, set
\[
a_y:=y^2+\delta_{\mathrm{ph}}^2+\varepsilon^2,
\qquad
b_y:=y\,\delta_{\mathrm{ph}},
\qquad
r_y:=\frac{a_y-\sqrt{a_y^2-4b_y^2}}{2b_y}.
\]
Then
\[
a_y-2b_y=(y-\delta_{\mathrm{ph}})^2+\varepsilon^2>0,
\qquad
0<r_y<1.
\]
If $d=d_{C_6}(i,j)$ denotes the cyclic graph distance on the hexagon, the regularized Green
kernel satisfies
\[
(\mathcal Y_y^{(\varepsilon)})_{ij}
=
\frac{r_y^{d}+r_y^{6-d}}
{\sqrt{a_y^2-4b_y^2}\,(1-r_y^6)}.
\]
Hence every matrix entry of $\mathcal Y_y^{(\varepsilon)}$ is strictly positive.
\end{proposition}

\begin{proof}[Proof sketch]
For the admissible cusps one has $y>0$ and $\delta_{\mathrm{ph}}\in(1/3,2/3)$, so $b_y>0$. The
displayed identity for $a_y-2b_y$ implies $a_y>2b_y$, hence $a_y^2-4b_y^2>0$ and the smaller
root $r_y$ lies strictly between $0$ and $1$. The operator
\[
\mathcal Y_y^{(\varepsilon)}=(a_y\mathbf 1-b_y(U_6+U_6^\dagger))^{-1}
\]
is circulant on the six-cycle, so its entries depend only on the cyclic distance
$d=d_{C_6}(i,j)$. Solving the corresponding Green recursion on $C_6$ gives the displayed closed
form, and every factor in that formula is strictly positive.
\end{proof}

\begin{corollary}[Uniform positive lower bound on admissible local transport]
\label{cor:uniform-positive-c6-green}
On the admissible cusp set
\[
y\in\left\{1,\frac13,\frac23\right\}
\]
and on the closed branch $\varepsilon=0$, all matrix entries of $\mathcal Y_y^{(0)}$ are bounded
below by a uniform constant
\[
g_{\mathrm{hex}}>0.
\]
One may take
\[
g_{\mathrm{hex}}\approx \num{0.6738706016}.
\]
\end{corollary}

\begin{proof}
There are only finitely many admissible cusp values and only finitely many matrix entries on
$C_6$. By \TFPTxref{prop:strict-positive-c6-green}, each of these entries is strictly positive at
$\varepsilon=0$, so their minimum defines a positive constant $g_{\mathrm{hex}}$.
\end{proof}

\begin{remark}[Carrier rank echo in the transport polynomial]\label{rem:transport-rank-echo}
On the minimal carrier branch $\gcar=5$, the hexagonal cosine polynomial of \Cref{eq:c6-cosine-polynomial} can be rewritten as
\[
(\gcar-1)C^4-\gcar C^2+\mathbf 1=0.
\]
Thus the transport closure polynomial carries the same pair $(\gcar-1,\gcar)$ that governs the hard carrier packet. Equivalently,
\[
4C^4-5C^2+\mathbf 1=(4C^2-\mathbf 1)(C^2-\mathbf 1)=(2C-\mathbf 1)(2C+\mathbf 1)(C-\mathbf 1)(C+\mathbf 1).
\]
So the distinguished value $C=1/2$ reappears on the transport side, mirroring the positive carrier root $Y=1/2$ of the factorized carrier polynomial in Section~3.3. At present this should be read as an exact compression identity on the retained hexagonal branch, not as a derivation of the transport sector from the carrier theorem.
\end{remark}

Absolute masses are attached to admissible path lengths,
\begin{equation}
\label{eq:path-length-yukawa}
y_{f,j}^{\mathrm{UV}}=\lambdaY^{L_{f,j}}\Lambda_{f,j},
\qquad
\lambdaY=\sqrt{\phiz(1-\phiz)}.
\end{equation}
The primary matrix-level statement is the hard holonomy form
\begin{equation}
\label{eq:yukawa-matrix-factorization}
Y_f=D_{L,f}U_f^\star D_{R,f},
\qquad
U_f^\star\in SU(3)_F,
\end{equation}
with $D_{L,f}$ and $D_{R,f}$ the diagonal word-length factors and $U_f^\star$ fixed by the
canonical family holonomy rather than treated as a residual free transport matrix. Equivalently,
\begin{equation}
\label{eq:yukawa-holonomy-matrix}
(Y_f)_{ij}
=
\lambdaY^{L^L_{f,i}+L^R_{f,j}}
\left\langle e_i,\mathrm{Hol}_{\Gamma_{ij}^{\min}}(\nabla_F^\star)e_j\right\rangle.
\end{equation}
\begin{theorem}[Exact family-holonomy phase]
\label{prop:holonomy-triangle-phase}
\label{prop:exact-family-holonomy-phase}
Define the ordered closed family holonomy by
\[
\Pi_\triangle
:=
\mathrm{Hol}_{\Gamma_{12}^{\min}}(\nabla_F^\star)\,
\mathrm{Hol}_{\Gamma_{23}^{\min}}(\nabla_F^\star)\,
\mathrm{Hol}_{\Gamma_{31}^{\min}}(\nabla_F^\star),
\qquad
\delta_{\mathrm{CKM}}
:=
\arg\tr\Pi_\triangle.
\]
This exact holonomy phase is the closed-branch phase quantity. On the
small transport-area branch one has
\begin{equation}
\label{eq:ckm-holonomy-phase}
\delta_{\mathrm{CKM}}
=
\frac{\pi}{3}+3\lambda_C^2+O(\lambda_C^4).
\end{equation}
\end{theorem}

\begin{proof}[Proof sketch]
The family holonomy is already encoded in \Cref{eq:yukawa-holonomy-matrix}. Closing the
minimal triangle isolates the gauge-invariant phase information carried by the family
connection. The $\mathbb Z_6$ transport grammar fixes the base phase at $\pi/3$, while the
first nontrivial nonabelian area correction is quadratic in the Cabibbo seed. Keeping only the
leading admissible area term gives the displayed asymptotic expansion of the exact phase.
\end{proof}

\TFPTxref{eq:path-length-yukawa} is the diagonal or singular-value reduction of
\Cref{eq:yukawa-holonomy-matrix}, not the conceptual starting point of the transport sector.
After the geometric electroweak closure fixes $\vgeo$, the intrinsic mass rows are emitted through
\begin{equation}
\label{eq:mass-bridge}
\hat m_{f,j}:=\frac{y_{f,j}\vgeo}{\sqrt2}.
\end{equation}
These hatted quantities are the intrinsic masses on the retained branch. Comparison-surface masses
are obtained only after the declared comparison map:
\[
m_{f,j}^{\mathrm{obs}}:=\mathcal R_{\SM}(\hat m_{f,j}).
\]
The role of the transport matrix elements $\Lambda_{f,j}$ must be kept explicit. They are bounded $O(1)$ admissible amplitudes, not unrestricted fit knobs. The RG bookkeeping behind direct $\vgeo$ insertions is part of the declared comparison map discussed later in \TFPTxref{sec:rg-map-compiler-outputs}.

\begin{remark}[Singlet-projected cusp lift to the transport hexagon]
The transport cusps do not come from the raw two-point carrier spectrum alone. The carrier itself has
\[
\mathrm{Spec}(Y|_E)=\left\{-\frac13,\frac12\right\}.
\]
After exterior lift to the one-family packet,
\[
\mathrm{Spec}(Y|_{S^+})
=
\left\{0,-\frac23,\frac16,1,-\frac12,\frac13\right\},
\]
with multiplicities displayed in \TFPTxref{tab:one-family}. Restricting to the singlet channels and taking absolute values yields
\[
\{|Y(e^c)|,\ |Y(u^c)|,\ |Y(d^c)|\}
=
\left\{1,\frac23,\frac13\right\},
\]
which are exactly the admissible cusp values of the transport hexagon. In this sense the transport sector reads the singlet-projected, family-lifted hypercharge spectrum through the resolvent
\[
D_y=y\mathbf 1-\delta_{ph}U_6
\]
on $C_6$. The transport grammar therefore lives on the same rational grid as the carrier packet, but only after exterior lift and singlet projection.
\end{remark}

On the closed branch, the retained diagonal word-length packets are
\begin{equation}
\label{eq:minimal-word-lengths}
\begin{gathered}
L_u^{\mathrm{diag}}=(L_u,L_c,L_t)=(7,3,0),\\
L_d^{\mathrm{diag}}=(L_d,L_s,L_b)=(7,5,2),\\
L_e^{\mathrm{diag}}=(L_e,L_\mu,L_\tau)=(8,5,3).
\end{gathered}
\end{equation}
These are not free parameters; they are the retained diagonal word-length packets forced by the
residue data and the closed-branch sum rule, with transport matrix elements in
\TFPTxref{tab:yukawa-residuals} remaining in the bounded $O(1)$ range. Once the hierarchy ordering is
fixed, the admissible graph distance is determined up to full $6$-step windings by
$\lambdaY\approx\num{0.2244}$.

\begin{remark}[Residue classes and full windings on the hexagon]
The exact resolvent formula \Cref{eq:finite-hexagonal-resolvent} makes the combinatorics behind \Cref{eq:minimal-word-lengths} transparent: every path length splits into a residue class on the base hexagon plus an optional full winding by $6$. Accordingly,
\[
L_t=0,\qquad
L_b=2,\qquad
L_\tau=L_c=3,\qquad
L_s=L_\mu=5
\]
live on the first hexagon, while
\[
\;L_u=L_d=1+6,
\qquad
L_e=2+6
\]
carry one additional full turn. Since
\[
\lambda_Y^6=\bigl(\phiz(1-\phiz)\bigr)^3\approx \num{1.28e-4},
\]
each extra winding produces the severe suppression needed for the first generation. This is not by itself a full mass theorem, but it is a structural reading of why the retained word-length pattern is so uneven.
\end{remark}
